\documentclass[11pt,aspectratio=169]{beamer}

% --------------------
% Theme & Packages
% --------------------
% \usetheme{Boadilla}
\usetheme{Madrid}
\usecolortheme{default}
\usefonttheme{professionalfonts}  % 数学公式字体
% \setbeamercolor{block title}{fg = white,bg = blue}
\usepackage{ctex}
\usepackage{amsmath, amssymb}
\usepackage{listings}
\usepackage{xcolor}

% --------------------
% Code Style
% --------------------
\lstset{
    language=C++,
    basicstyle=\ttfamily\small,
    keywordstyle=\color{blue},
    stringstyle=\color{red},
    commentstyle=\color{gray},
    numbers=left,
    numberstyle=\tiny,
    stepnumber=1,
    breaklines=true,
    frame=single
}

% --------------------
% Title Info
% --------------------
\title{绪论}
\author{qwaszx}
\date{\today}

% ====================
% Document
% ====================
\begin{document}

% --------------------
\begin{frame}
  \titlepage
\end{frame}

% --------------------
\begin{frame}{今天的内容}
\begin{itemize}
    \item 为什么学 OI / 学 OI 有什么好处？
    \item 怎么学 OI？
    \item 对本系列课程的介绍
\end{itemize}
\end{frame}

% =========================================================
% Part 1
% =========================================================
\section{为什么要学 OI}

\begin{frame}{为什么学了 OI}
    你为什么学了 OI？

    \pause
    我为什么学了 OI。
\end{frame}

\begin{frame}{学 OI 的好处}
\begin{itemize}
    \item 提升逻辑思维能力
    \item 算法思维
    \begin{itemize}
        \item 如何描述解决问题的步骤
        \item 对计算的理解
    \end{itemize}
    \item 计算机和离散数学知识
    \item OI 社群
    \item 更多地玩电脑
    \item 升学与就业
\end{itemize}
\end{frame}

\begin{frame}{升学与就业}
升学方面：
    \begin{itemize}
        \item 金牌可以保送清华姚班/北大图灵班
        \item 银牌可以破格入围强基计划
        \item 省一对强基计划和各种自主招生有一定帮助 
    \end{itemize}
就业方面：
    \begin{itemize}
        \item OI/XCPC 竞赛获奖对于计算机相关行业有一定帮助
        \item 可以继续做教培，相对普通文化课报酬更高
        \item 对口 TCS（理论计算机科学）
    \end{itemize}
\end{frame}

\begin{frame}{好处说完了，那么坏处呢}
    \begin{itemize}
        \item 机会成本：选择 OI 需要放弃一些其他东西
        \item 更多地使用计算机和互联网，可能会损害身心健康。
        \item 对于升学和就业的性价比不算很高
        \begin{itemize}
            \item 目前，银牌相比省一的优势不大，但是困难得多
            \item 集训队每年只有 50 人，竞争激烈
        \end{itemize}
    \end{itemize}
\end{frame}

\begin{frame}{总结}
    初中阶段，适合学文化课有余力并且学 OI 学得开心的同学学习。

    高中阶段，功利角度建议根据自身成绩谨慎选择（高一进省队 -> 高二进集训队）

    兴趣是第一位的！
\end{frame}
% =========================================================
% Part 2
% =========================================================
\section{怎么学 OI}

\begin{frame}{OI 比赛的级别}
    \begin{enumerate}
        \item CSP-J
        \item CSP-S
        \item NOIP
        \item 省选
        \item NOI
        \item CTS
        \item IOI
    \end{enumerate}
    其中 CSP-J 到 CSP-S 属于前期，CSP-S 到 NOIP 属于中期，省选到 NOI 属于后期，再往后是大后期
\end{frame}

\begin{frame}{做题的过程}
\begin{enumerate}
    \item \textbf{读题}
    \item \textbf{分析性质}
    \item \textbf{抽象问题}
    \item \textbf{使用/改编/设计算法来处理}
    \item \textbf{实现}
\end{enumerate}

前期最重要的是 1 和 5，中期最重要的是 3 和 4，后期和大后期最重要的是 2 和 4

学习的过程中，每个部分都在逐步提高
\end{frame}

\begin{frame}{读题与实现}
    抽象/封装：忽略某个东西的内部结构，只关注其输入和输出

    在做题中练习

    多看看其他人的写法
\end{frame}

\begin{frame}{抽象问题，调用算法}
    我们会学习很多的算法。学习的过程中，
    \begin{itemize}
        \item 关注算法整体的功能，熟知一个问题可以被什么算法解决，一个算法可以解决什么问题，有哪些变种和扩展
        \item 关注算法的每一步，熟知其原理、作用，多思考可以如何扩展
        \item 关注证明
    \end{itemize}

    在做题中练习

    多总结
\end{frame}

\begin{frame}{分析性质，设计算法}
    也就是所谓的人类智慧。一般来说，很难！

    多总结题目和算法的模型。相似的模型往往具有相似的性质。

    多思考题解/算法的每一步是如何想到的。

    尝试建立不同概念和不同算法之间的联系，抽象其共性
\end{frame}

\begin{frame}{训练与学习资源}
    OJ/题库：洛谷、QOJ、LOJ、UOJ；Codeforces、Atcoder

    比赛：Codeforces、Atcoder

    学习资源：oiwiki、各种优质博客（如 liu\_cheng\_ao 推荐的、command\_block、Elegia）

    各种 AI

    OI 社群与搜索引擎。不懂就问/搜！
\end{frame}

\begin{frame}{训练的方法}
    \begin{itemize}
        \item 学习知识点
        \item 刷知识点对应的题
        \item 对知识点进行归纳总结
        \item 通过模拟赛进行不同知识点的综合练习/复习
    \end{itemize}
\end{frame}

\begin{frame}{刷题的方法}
    先自行思考，思考到自己一直没有任何想法为止，然后看题解。这个时间长短是不定的

    看题解的过程中也要进行思考：作者说的有没有道理，作者是怎么想到的，我能不能自行补全剩下的部分，哪些部分还可以优化

    想出做法/看懂题解后进行实现，这个过程中不要参考任何现成代码

    做完题后看看题解区有没有不同做法，看看其他人的代码都是怎么写的，想想自己的代码哪些地方能更快/更好写

    看看题目是否是\textbf{经典问题}，是则需要在归纳总结时\textbf{记录}
\end{frame}

\begin{frame}{经典问题}
    有一些问题是经典的：在拆掉包装后，其技巧/模型非常泛用。每做完一道题都应该思考这道题是不是这种问题。换句话说，做题后要努力将题目的模型/技巧总结出来。
\end{frame}

\begin{frame}{刷什么题}
    学习资源中提到的题、自行根据知识点搜索相关题目（如洛谷题单）

    重要的是刷题前中后的\textbf{思考}而不是刷题本身。所以，非常同质化的题不用多刷，几道就够了（再多的可以口胡）

    刷题可以形成题目<->模型<->算法的条件反射，或者说直觉。
\end{frame}

\begin{frame}{训练的理念}
    循序渐进，避免一直学太简单/太难的东西

    青少年的智力、理解能力会随着年龄提高

    熟练掌握 OI 大部分知识点通常需要且只需要几千个小时，相当于 1 到 2 年*8h（高中脱产/初中或者小学就开始学） 
\end{frame}

\begin{frame}{OI 之外}
    对初中生来说，我非常建议学好文化课，义务教育是必要的。不过可以多和家长老师沟通，放弃重复性过强的文化课练习。OI 以兴趣为主，无需强求。

    OI 涉及大量的思考，所以充足睡眠非常重要，建议想办法得到充足睡眠。最起码的原则是不能一天比一天困。

    多锻炼。一方面大学也有各种体育测试（如 3km），另一方面人生是场马拉松，健康地活着才是成功

    人际交流，线上线下均可
\end{frame}

% =========================================================
% Part 3
% =========================================================
\section{语法基础 Revisit}

\begin{frame}{代码能力}
    可以用代码实现想法

    知道一个东西能不能写、大概是怎么写的、好不好写
    
    不出 bug 地尽量快地写出来

    学习更优秀的写法（更好写/更快）
\end{frame}

\begin{frame}{先想再写 vs 边写边想}
    我的习惯：在写某个部分时，对这个部分的整体结构必须提前想好，细节可以边写边想

    抽象/封装的代码设计理念
\end{frame}

\begin{frame}{目前你需要非常熟练的东西}
    \begin{itemize}
        \item 顺序、分支、循环结构的程序设计
        \item 各种变量类型、结构体、数组
        \item 函数
        \item 简单的模拟、枚举
    \end{itemize}

    如果不熟练，可以先多做做橙题
\end{frame}

\section{对本课程的介绍}
\begin{frame}{培养目标}
    在课程开始前，我们希望大家熟练掌握基本的 C++ 和算法知识，可以熟练切掉橙题；此外，最好基本学过一遍 CSP-S 范围内的算法。

    在本期课程中，我们会较为深入地学习各种 CSP-S 级别的算法，并且在做题的过程中看到各种相关的技巧和模型，期望让大家对这些算法有较为深入的理解，遇到对应的问题时可以轻松解决。

    如果有机会的话，在后续的课程中，我们会更深入地学习这些算法相关的内容，期望大家在这部分内容上达到 NOIP/省选/NOI 的水平（注意：这三个东西的难度其实没有显著区别）。不过由于这部分内容密度较大，学习起来可能较为困难。
\end{frame}

\begin{frame}{课程节奏}
    每周一次课 + 每周题单（10 道题左右） + 一次模拟赛

    课下希望大家多多练习，大概每天做两道题左右即可
\end{frame}
% =========================================================
% Part 5
% =========================================================
\section{总结}

\begin{frame}{总结}
\begin{itemize}
    \item 对 OI 概况、利弊的了解
    \item OI 的一些学习方法
    \item 本课程的介绍
\end{itemize}
\end{frame}

\begin{frame}{最后}
希望大家学得开心

\pause
\centerline{\Large 谢谢大家！}
\end{frame}

\end{document}
